\subsection{Sprint 0}

Na Sprint 0, meu foco esteve concentrado na estruturação inicial do produto e na organização do trabalho do time. Comecei atuando diretamente na criação de \textit{user stories}, principalmente as relacionadas ao Painel Admin, e fiquei envolvido de forma completa na definição das 7 histórias dessa frente. Esse trabalho foi importante para deixar mais claro o que realmente precisava ser construído no início do projeto e para reduzir ambiguidades na hora de transformar as demandas em tarefas técnicas.

Além das histórias, também criei tarefas voltadas ao gerenciamento de edificações tanto para \textit{frontend} quanto para \textit{backend}. Em paralelo, comecei a mapear tarefas relacionadas a gerenciamento de imagens e idiomas. Essas últimas ainda ficaram em fase de refinamento, porque preferimos amadurecer melhor os critérios antes de levar para execução direta. Mesmo assim, esse levantamento inicial me ajudou a antecipar pontos que poderiam gerar retrabalho mais adiante.

Outra atividade relevante dessa sprint foi a preparação para a reunião com os clientes. Organizei perguntas para conduzir melhor a conversa e, junto com os outros AGES IV, conseguimos obter respostas objetivas sobre o que era prioridade no produto. Esse momento foi decisivo para alinhar expectativas e para validar o direcionamento das histórias do admin, já que boa parte das decisões de escopo dependia desse retorno. A reunião também serviu para fortalecer a conexão com os \textit{stakeholders}, acabei sendo o responsável por conduzir esta conversa, sendo essa a minha primeira experiência com esse tipo de contato direto, o que foi bastante enriquecedor para entender melhor as necessidades do projeto e para desenvolver habilidades de comunicação e alinhamento e criar uma boa relação com os clientes.

Como dificuldade, no começo eu me senti perdido para definir com clareza quem ficaria responsável por cada parte entre os AGES IV, e isso acabou atrasando um pouco o andamento inicial. A principal lição que levo dessa sprint é a necessidade de distribuir responsabilidades de forma mais assertiva desde o início, com escopos melhor definidos para cada pessoa. Quando a divisão fica clara, o time ganha ritmo mais rápido e as entregas tendem a ficar mais consistentes.

No mais, essa Sprint foi tranquila em termos de execução e tudo parecia ocorrer bem, embora ainda estivessemos um pouco perdidos, o que é natural para um projeto que está começando. O importante é que aula por aula conseguíamos nos conectar melhor e entender as necessidades do projeto, assim como de cada aluno, que nos deixou confiante para o continuamento da jornada.
